% S04-M6 (package abhijit/s04-math): the eta-principle (commutative algebra), maths only.
% Labels: prop:s04-etaprinciple.
% Reference statement: the S06-L17 row of plan/implementation/stage-06.md, section 1 (agreed single reference for this text and the
% docstring of Langlandscape.CommAlg.EtaPrinciple.congruenceIdeal_eq_span_pairing).
\section{The \texorpdfstring{$\eta$}{eta}-principle}

This section proves the commutative-algebra statement behind Thm.~\ref{thm:s04-etaB} (the paragraph after it in the stage tex): a free
rank-one module over a finite free algebra, with a perfect equivariant pairing, computes the congruence ideal as the self-pairing of a
generator of the eigenline. It is Hsieh's Remark~7.3 for $d=1$
(\lit{hsieh\_2017\_hida-families-and-p-adic-triple-product-l-functions\_arxiv}{Remark~7.3, p.~67}), without his assumptions that the base
ring is a domain and the algebra is reduced. Gorenstein-ness is not assumed; it follows from the pairing (part (a)).

Conventions, as in stage~06 (\texttt{if:s06-lemD}, Lean namespace \pathtt{Langlandscape.CommAlg}): for a commutative ring $A$ and a
commutative $A$-algebra $S$, a linear form $\ell\in\Hom_A(S,A)$ is a \emph{trace} (and $S$ is \emph{Gorenstein with trace $\ell$},
$\mathrm{IsGorensteinWith}\,A\,S\,\ell$) if $S$ is a finite free $A$-module and $s\mapsto\ell(s\,\cdot)$ is a bijection
$S\to\Hom_A(S,A)$. For an ideal $J\subseteq S$, $\Ann_S J=\{s:sJ=0\}$. For an $A$-algebra map $\mu:S\to A$, the \emph{congruence ideal}
is $\mathrm{congruenceIdeal}(\mu):=\mu(\Ann_S\ker\mu)\subseteq A$ (an ideal, as $\mu$ is surjective); this is $\eta(S,\mu)$ of
seed \S2.6 and Hsieh's $C(\lambda)$.

\begin{proposition}[$\eta$-principle; \pathtt{thm:s04-etaB}, \texttt{lem:s06-S}]\label{prop:s04-etaprinciple}
Let $A$ be a commutative ring and $S$ a commutative $A$-algebra that is finite free as an $A$-module. Let $M$ be an $S$-module with an
$S$-linear isomorphism $e:M\xrightarrow{\sim}S$, and let $B:M\times M\to A$ be an $A$-bilinear form such that
\begin{itemize}
\item $B$ is \emph{perfect}: the $A$-linear map $M\to\Hom_A(M,A)$, $m\mapsto B(m,\cdot)$, is bijective;
\item $B$ is \emph{$S$-equivariant}: $B(sm,n)=B(m,sn)$ for all $s\in S$, $m,n\in M$.
\end{itemize}
Let $\mu:S\to A$ be a homomorphism of $A$-algebras, and let $m_0\in M$ generate the eigenline
\[M[\mu]:=\{m\in M:\ sm=\mu(s)m\ \text{for all }s\in S\},\qquad\text{i.e.\ }M[\mu]=A\,m_0.\]
Then:
\begin{enumerate}[label=(\alph*)]
\item $S$ is Gorenstein with trace $\ell(s):=B(m_1,s\,m_1)$, where $m_1:=e^{-1}(1)$ (\lean{Langlandscape.CommAlg.EtaPrinciple.isGorensteinWith\_of\_pairing});
\item $\mathrm{congruenceIdeal}(\mu)=\big(B(m_0,m_0)\big)$, the principal ideal of $A$ generated by $B(m_0,m_0)$
  (\lean{Langlandscape.CommAlg.EtaPrinciple.congruenceIdeal\_eq\_span\_pairing}).
\end{enumerate}
\stproved{} below. \stformal{} (S06-L17, \pathtt{CommAlg/EtaPrinciple.lean}), in full generality of the statement and without registered hypotheses: \lean{Langlandscape.CommAlg.EtaPrinciple.isGorensteinWith\_of\_pairing} for (a) and \lean{Langlandscape.CommAlg.EtaPrinciple.congruenceIdeal\_eq\_span\_pairing} for (b); the trace in Lean is \texttt{traceOfPairing}, $\ell=B(e^{-1}(1),e^{-1}(\cdot))$.
\end{proposition}
\noindent The hypotheses are exactly those of the S06-L17 row of the stage-06 plan: $S$ finite free over $A$; $e:M\cong S$ as
$S$-modules; $B:M\to\Hom_A(M,A)$ bijective; $B(s\cdot m,n)=B(m,s\cdot n)$; $\mu:S\to A$ an $A$-algebra map; $m_0$ generating the
eigenline $\{m:\forall s,\ s\cdot m=\mu(s)\cdot m\}$. No locality, noetherianity, domain or reducedness hypothesis is needed.

\begin{proof}
The element $\theta$ below is the one of Lemma~S of stage~06 (\texttt{lem:s06-S}); we give the argument in full.

\emph{(a)} For $s,t\in S$, $e^{-1}(s)=s\,m_1$, and by equivariance
\[B(s\,m_1,t\,m_1)=B(m_1,st\,m_1)=\ell(st).\]
Hence the map $S\to\Hom_A(S,A)$, $s\mapsto\ell(s\,\cdot)$, is the composite
\[S\xrightarrow{\ e^{-1}\ }M\xrightarrow{\ m\mapsto B(m,\cdot)\ }\Hom_A(M,A)\xrightarrow{\ \varphi\mapsto\varphi\circ e^{-1}\ }\Hom_A(S,A)\]
of three bijections, so it is bijective. With $S$ finite free over $A$ this is $\mathrm{IsGorensteinWith}\,A\,S\,\ell$. (In passing:
$B(sm_1,tm_1)=\ell(st)=B(tm_1,sm_1)$, so $B$ is symmetric.)

\emph{(b), Step 1: the annihilator.} As $\ell(\cdot\,\cdot)$ is bijective, there is a unique $\theta\in S$ with $\ell(\theta x)=\mu(x)$ for all
$x\in S$. Note $\ker\mu=\{x-\mu(x)\cdot1:x\in S\}$ (as $\mu$ is $A$-linear with $\mu(1)=1$). We claim $\Ann_S\ker\mu=A\theta$ and
$a\theta=0\Rightarrow a=0$ for $a\in A$.
(i) $\theta\in\Ann_S\ker\mu$: for $y\in\ker\mu$ and all $x$, $\ell(\theta y\,x)=\mu(yx)=\mu(y)\mu(x)=0$, so $\theta y=0$ by injectivity
of $s\mapsto\ell(s\,\cdot)$.
(ii) If $z\in\Ann_S\ker\mu$, then $z(x-\mu(x))=0$, i.e.\ $zx=\mu(x)z$, for all $x$; so $\ell(zx)=\mu(x)\ell(z)=\ell(\ell(z)\theta\,x)$ for all
$x$, whence $z=\ell(z)\theta$ by injectivity.
(iii) If $a\theta=0$ then $0=\ell(a\theta\cdot1)=a\mu(1)=a$.
So $\Ann_S\ker\mu=A\theta$ is free of rank one over $A$, and $\mathrm{congruenceIdeal}(\mu)=\mu(A\theta)=(\mu(\theta))$.

\emph{Step 2: the eigenline.} For $m\in M$ we have $sm=\mu(s)m$ for all $s$ iff $(s-\mu(s))e(m)=0$ for all $s$ iff
$e(m)\in\Ann_S\ker\mu$. So $M[\mu]=e^{-1}(A\theta)=A\,\theta m_1$, and $a\mapsto a\,\theta m_1$ is injective by Step~1(iii).

\emph{Step 3: generators.} Since $A m_0=M[\mu]=A\,\theta m_1$, write $m_0=a\,\theta m_1$ and $\theta m_1=b\,m_0$ with $a,b\in A$. Then
$\theta m_1=ab\,\theta m_1$, so $(1-ab)\theta m_1=0$ and $ab=1$ by Step~2: $a$ is a unit.

\emph{Step 4: the self-pairing.} $B(\theta m_1,\theta m_1)=\ell(\theta\cdot\theta)=\mu(\theta)$, so
$B(m_0,m_0)=a^2B(\theta m_1,\theta m_1)=a^2\mu(\theta)$. As $a$ is a unit, $(B(m_0,m_0))=(\mu(\theta))=\mathrm{congruenceIdeal}(\mu)$.
\end{proof}

\noindent\emph{Remarks.} (1) A generator $m_0$ always exists (Step~2: $\theta m_1$ is one), and any two differ by a unit of $A$; so
$B(m_0,m_0)$ is well defined up to the square of a unit, and the ideal $(B(m_0,m_0))$ is independent of all choices.
(2) Step~1 is Lemma~S of stage~06 (\texttt{lem:s06-S}, without base change) for the trace $\ell$ of (a); in Lean the proof of (b) can
use \pathtt{LemmaS.congruenceIdeal\_eq\_span} with $\ell$ from (a), which is why S06-L17 is stacked on S06-L2.
(3) For Hsieh's application (Thm.~\ref{thm:s04-etaB}) one takes $A=\mathbb I$, $S=\T^D(N,\chi,\mathbb I)_\m$,
$M=eS^D(N,\chi,\mathbb I)_\m$, $B=B_N$, $\mu=\lambda^D_F$, $m_0=F^D$; the freeness $M\cong S$ is Thm.~\ref{thm:s04-Gor} (cited), and the
perfectness and equivariance of $B_N$ are [Hs, Def.~4.3] (cited).
(4) Hsieh's Remark~7.3 treats free $H$ of rank $d$ with $C(\lambda)^d=(\det\langle e_i,e_j\rangle)$; the case $d=1$ is (b). The
general case is not needed in stages 04--06.
